%% LyX 2.4.4 created this file.  For more info, see https://www.lyx.org/.
%% Do not edit unless you really know what you are doing.
\documentclass[english,handout]{beamer}
\usepackage{mathptmx}
\usepackage[T1]{fontenc}
\usepackage[latin9]{inputenc}
\usepackage{units}
\usepackage{amssymb}
\usepackage{graphicx}

\makeatletter
%%%%%%%%%%%%%%%%%%%%%%%%%%%%%% Textclass specific LaTeX commands.
% this default might be overridden by plain title style
\newcommand\makebeamertitle{\frame{\maketitle}}%
% (ERT) argument for the TOC
\AtBeginDocument{%
  \let\origtableofcontents=\tableofcontents
  \def\tableofcontents{\@ifnextchar[{\origtableofcontents}{\gobbletableofcontents}}
  \def\gobbletableofcontents#1{\origtableofcontents}
}

%%%%%%%%%%%%%%%%%%%%%%%%%%%%%% User specified LaTeX commands.
\usetheme{Warsaw}
% or ...

\setbeamercovered{transparent}
% or whatever (possibly just delete it)
\setbeamercovered{invisible} %makes overlays invisible


%gets rid of navigation symbols
\setbeamertemplate{navigation symbols}{}


\usepackage{pgfpages}
\pgfpagesuselayout{4 on 1}[a4paper, landscape, border shrink=7mm]
\pgfpageslogicalpageoptions{1}{border code=\pgfusepath{stroke}}
\pgfpageslogicalpageoptions{2}{border code=\pgfusepath{stroke}}
\pgfpageslogicalpageoptions{3}{border code=\pgfusepath{stroke}}
\pgfpageslogicalpageoptions{4}{border code=\pgfusepath{stroke}}

\makeatother

\usepackage{babel}
\begin{document}
\title[Short Title]{9.5: The Ratio, Root, and Comparision Tests.}
\subtitle{Calculus III | MTH 331~~|~~Section A~~|~~Fall 2012}
\author{Dr.\,James Ashe}
\titlegraphic{\includegraphics[height=1.5cm]{../../jcsu_bull.jpg}}
\institute{Johnson C. Smith University}
\date{{}}

\makebeamertitle

%\beamerdefaultoverlayspecification{<+->}
\begin{frame}{The Ratio Test}

\begin{block}{Theorem. The Ratio Test}

Given a series $\sum a_{k}$, let $r={\displaystyle \lim_{k\rightarrow\infty}\left|\frac{a_{k+1}}{a_{k}}\right|}$.

\begin{enumerate}
\item If $0\leq r<1$ the series converges.
\item If $1<r$ then the series diverges.
\item If $r=1$ the test is inconclusive.Use
\end{enumerate}
\end{block}
\begin{itemize}
\item <2->[]Use the Ratio Test to determine whether the following series
converge.
\item <3->[ex 9.]${\displaystyle {\displaystyle \sum_{k=1}^{\infty}\frac{1}{k!}}}$. 
\item <4->[]${\displaystyle \lim_{k\rightarrow\infty}\frac{a_{k+1}}{a_{k}}=\lim_{k\rightarrow\infty}\frac{\nicefrac{1}{(k+1)!}}{\nicefrac{1}{k!}}=\lim_{k\rightarrow\infty}\frac{k!}{(k+1)!}}$\uncover<5->{\\
$={\displaystyle \lim_{k\rightarrow\infty}\frac{1\cdot2\cdot3\cdots k}{1\cdot2\cdot3\cdots k\cdot(k+1)}=\lim_{k\rightarrow\infty}\frac{1}{(k+1)}=0}$.}\uncover<6->{
$\Rightarrow\sum_{k=1}^{\infty}\frac{1}{k!}$ converges.}
\end{itemize}
\end{frame}

\begin{frame}{The Ratio Test}

Why does the Ratio Test work?
\begin{itemize}
\item <2->As $k$ gets large, $r={\displaystyle \lim_{k\rightarrow\infty}\left|\frac{a_{k+1}}{a_{k}}\right|}\Rightarrow r\cdot a_{k}\approx a_{k+1}$.
\item <3->It is the ``tail'' of the series that determines its end behavior.
\item <4->So then $a_{k}+a_{k+1}+a_{k+2}+\dots\approx a_{k}+r^{1}\cdot a_{k+1}+r^{2}\cdot a_{k+2}+\dots\approx$$a_{k}(1+r+r^{2}+\dots)$
behaves geometrically;
\item <5->[]geometric series converge when $r<1$ and diverge when $r\geq1$. 
\end{itemize}
\end{frame}

\begin{frame}{The Ratio Test}

\begin{itemize}
\item [ex.]${\displaystyle {\displaystyle \sum_{n=1}^{\infty}\frac{9^{n}}{(-2)^{n+1}n}}}$. 
\item <2->Note: $\frac{a_{n+1}}{a_{n}}=a_{n+1}(\frac{1}{a_{n}})$
\item <2->[]${\displaystyle \lim_{n\rightarrow\infty}}\left[a_{n+1}\left(\frac{1}{a_{n}}\right)\right]={\displaystyle \lim_{n\rightarrow\infty}\left[\frac{9^{n+1}}{(-2)^{n+2}(n+1)}\frac{(-2)^{n+1}n}{9^{n}}\right]}$\textrm{$=\lim_{n\rightarrow\infty}\frac{9n}{(-2)^{1}(n+1)}=-\frac{9}{2}{\displaystyle \lim_{n\rightarrow\infty}\frac{n}{(n+1)}}$\uncover<3->{$=-\frac{9}{2}$}}
\item <4->[], and $\left|-\frac{9}{2}\right|>1\Rightarrow$\uncover<5->{$\sum_{n=1}^{\infty}\frac{9^{k}}{(-2)^{n+1}n}$
diverges.}
\end{itemize}
\end{frame}

\begin{frame}{The Root Test}

\begin{block}{Theorem: The Root Test}

Let $\sum a_{k}$ be an infinite series with \emph{nonnegative }terms
and let $\rho={\displaystyle \lim_{k\rightarrow\infty}\sqrt[k]{a_{k}}}$,

\begin{enumerate}
\item <2->If $0\leq\rho<1$, then the series converges.
\item <3->If $1<\rho$, then the series diverges.
\item <4->If $\rho=1$, the test is inconclusive.
\end{enumerate}
\end{block}
\begin{itemize}
\item [ex 20.]<5->Evaluate${\displaystyle \sum_{k=1}^{\infty}\left(\frac{k+1}{2k}\right)^{k}}$.
\item []<6->$\sqrt[k]{\left(\frac{k+1}{2k}\right)^{k}}=\left(\frac{k+1}{2k}\right)$.
\uncover<7->{${\displaystyle \lim_{n\rightarrow\infty}\frac{n+1}{2n}}=$}\uncover<8->{${\displaystyle \lim_{k\rightarrow\infty}\frac{\nicefrac{n}{n}+\nicefrac{1}{n}}{2\nicefrac{n}{n}}}=\frac{1+0}{2}=\frac{1}{2}$.}
\end{itemize}
\end{frame}

\begin{frame}{The Root Test}

Why does the Root Test work?
\begin{itemize}
\item <2->As $k$ gets large, $\rho\approx\sqrt[k]{a_{k}}$ or $a_{k}\approx\rho^{k}$.
\item <3->Again, it is the ``tail'' of the series that determines its
end behavior.
\item <4->So then $a_{k}+a_{k+1}+a_{k+2}+\dots\approx\rho^{k}+\rho^{k+1}+\rho^{k+2}+\dots$
behaves geometrically;
\item <5->[]geometric series converge when $\rho<1$ and diverge when $r\geq\rho$. 
\end{itemize}
\end{frame}

\begin{frame}{The Comparison Test}

\begin{block}{Theorem. Comparison Test}

Let $\sum a_{k}$ and $\sum b_{k}$ be series with positive terms.

\begin{enumerate}
\item If $0<a_{k}\leq b_{k}$ and $\sum b_{k}$ converges then $\sum a_{k}$
also converges.
\item If $0<b_{k}\leq a_{k}$ and $\sum b_{k}$ diverges then $\sum a_{k}$
also diverges.
\end{enumerate}
\end{block}
\begin{itemize}
\item <2->[ex 31.]Evaluate ${\displaystyle \sum_{k=1}^{\infty}\frac{k}{k^{2}-\cos^{2}(k)}}$.
\item <3->[]$0\leq\cos^{2}(x)\leq1$
\item <3->[]$\frac{k}{k^{2}-\cos^{2}(k)}>\frac{k}{k^{2}}=\frac{1}{k}$
(making the denominator smaller$\Rightarrow$fraction gets bigger)
\item <4->[]${\displaystyle \sum_{k=1}^{\infty}\frac{1}{k}=\infty}$ (the
Harmonic Series) so then $\sum_{k=1}^{\infty}\frac{k}{k^{2}-\cos^{2}(k)}$
\uncover<5->{also diverges.}
\end{itemize}
\end{frame}

\begin{frame}{The Limit Comparison Test}

\begin{block}{Theorem. Limit Comparison Test}

Define $L={\displaystyle \lim_{k\rightarrow\infty}\frac{a_{k}}{b_{k}}}$
where $a_{n},b_{n}\geq0$ for all $n$. 

\begin{enumerate}
\item If $0<L<\infty$ then $\sum a_{k},\sum b_{k}$ either both converge
or diverge.
\item If $L=0$ and $\sum b_{k}$ converges then $\sum a_{k}$ converges.
\item If $L=\infty$ and $\sum a_{k}$ diverges then $\sum b_{k}$ diverges.
\end{enumerate}
\end{block}
\begin{itemize}
\item [ex.]<2->Evaluate ${\displaystyle \sum_{k=0}^{\infty}\frac{1}{3^{k}-k}}$.
\uncover<3->{guess: it converges}
\item []<4->Need a series to compare (want it to converge): \uncover<5->{${\displaystyle \sum_{k=0}^{\infty}}\frac{1}{3^{k}}$.}
\item []<6->${\displaystyle \lim_{k\rightarrow\infty}\frac{\nicefrac{1}{3^{k}}}{\nicefrac{1}{(3^{k}-k)}}={\displaystyle \lim_{k\rightarrow\infty}\left(1-\frac{k}{3^{k}}\right)=1}}$\uncover<7->{$\Rightarrow{\displaystyle \sum_{k=0}^{\infty}\frac{1}{3^{k}-k}}$
converges since \textrm{${\displaystyle \sum_{k=0}^{\infty}}\frac{1}{3^{k}}$
does.}}
\end{itemize}
\end{frame}

\begin{frame}{}

Homework so far

\textbf{9.1}: 21, 55, 60, 66, 68 

\textbf{9.2}: 4, 9, 18, 20, 29, 46, 48, 51-56, 58 

\textbf{9.3}: 12, 19, 23, 26, 47, 49 

\textbf{9.4}: 12, 14, 20, 24, 26, 29 

\textbf{9.5}: 9-14, 19-25, 27-28
\end{frame}

\begin{frame}{}

\begin{block}{Recall: Integration by Parts}

$\int u\,dv=uv-\int v\,du$
\end{block}
\begin{itemize}
\item <2->[ex.]Find $\int xe^{-x}\,dx$. \uncover<3->{Let $u=x$ so that,}\uncover<4->{
\[
\begin{array}{rcl}
u=x &  & dv=e^{-x}\,dx\\
du=x &  & v=-e^{-x}
\end{array}
\]
}\vspace{-10pt}
\item <5->[]$\int\underset{u}{\underbrace{x}}\cdot\underset{dv}{\underbrace{e^{-x}\,dx}}$\uncover<6->{$=\underset{u}{\underbrace{x}}\underset{v}{\underbrace{(-e^{-x})}}-\int\underset{v}{\underbrace{(-e^{-x})}}\,\underset{du}{\underbrace{dx}}$.}
\item <7->[]$=-xe^{-x}-e^{-x}+C$
\end{itemize}
\begin{block}<8->{$e^{x}$}
${\displaystyle \lim_{n\rightarrow\infty}\left(\frac{n+x}{n}\right)^{n}=e^{x}}$
\end{block}
\uncover<9->{In particular, ${\displaystyle \lim_{n\rightarrow\infty}\left(\frac{n}{n+1}\right)^{n}=e^{-1}}$.}
\begin{itemize}
\item <9->[]These facts come in useful on problems \textbf{9.4: }29, \textbf{9.5:
}12 and 14 (among others)
\end{itemize}
\end{frame}

\end{document}
